% !TEX root = ../Hauptdatei.tex
% ============================================================
% Kapitel: Bilder
% ============================================================

\chapter{Bildgrößen und Skalierung}
\label{sec:groesse}
\index{Bild!Skalierung}

\section{Bild mit fester Breite}

Ein Bild, das auf 60\% der Textbreite skaliert wird:

\begin{figure}[htbp]
    \centering
    \includegraphics[width=0.6\textwidth]{example-image}
    \caption{Ein Beispielbild mit 60\,\% Textbreite}
    \label{fig:breite60}
\end{figure}

\section{Bild mit fester Höhe}

Ein Bild, dessen Höhe festgelegt wird:

\begin{figure}[htbp]
    \centering
    \includegraphics[height=4cm]{example-image}
    \caption{Ein Beispielbild mit 4\,cm Höhe}
    \label{fig:hoehe4}
\end{figure}

\section{Bild mit Skalierungsfaktor}

Skalierung um den Faktor 0{,}4:

\begin{figure}[htbp]
    \centering
    \includegraphics[scale=0.4]{example-image}
    \caption{Ein Beispielbild mit Skalierungsfaktor 0{,}4}
    \label{fig:scale04}
\end{figure}

\section{Breite und Höhe gleichzeitig}

Wenn beide angegeben werden, wird das Bild proportional so skaliert,
dass es in den Rahmen passt:

\begin{figure}[htbp]
    \centering
    \includegraphics[width=0.4\textwidth, height=3cm, keepaspectratio]{example-image}
    \caption{Bild mit \texttt{keepaspectratio}\index{keepaspectratio} -- verzerrungsfrei}
    \label{fig:keepaspect}
\end{figure}

% ============================================================
\chapter{Bildtransformationen}
\label{sec:transformation}
\index{Bild!Transformation}
% ============================================================

\section{Gedrehtes Bild}
\index{Bild!Drehen}

Ein um 15\,Grad gedrehtes Bild:

\begin{figure}[htbp]
    \centering
    \includegraphics[width=0.5\textwidth, angle=15]{example-image}
    \caption{Ein um 15\textdegree\ gedrehtes Bild}
    \label{fig:gedreht}
\end{figure}

\section{Bild mit Rahmen}
\index{Bild!Rahmen}

Ein Bild mit einem Rahmen versehen:

\begin{figure}[htbp]
    \centering
    \fbox{\includegraphics[width=0.5\textwidth]{example-image}}
    \caption{Ein gerahmtes Beispielbild}
    \label{fig:rahmen}
\end{figure}

\section{Gespiegeltes Bild}
\index{Bild!Spiegeln}

Horizontal gespiegelt (\texttt{scale=-1}):

\begin{figure}[htbp]
    \centering
    \includegraphics[width=0.4\textwidth, scale=-1]{example-image}
    \caption{Horizontal gespiegeltes Bild}
    \label{fig:gespiegelt}
\end{figure}

% ============================================================
\chapter{Bildanordnung}
\label{sec:anordnung}
\index{Bild!Anordnung}
% ============================================================

\section{Mehrere Bilder nebeneinander}

Zwei Bilder mit Untertiteln nebeneinander:

\begin{figure}[htbp]
    \centering
    \begin{minipage}{0.45\textwidth}
        \centering
        \includegraphics[width=\textwidth]{example-image-a}
        \subcaption{Variante A}
        \label{fig:nebeneinander-a}
    \end{minipage}
    \hfill
    \begin{minipage}{0.45\textwidth}
        \centering
        \includegraphics[width=\textwidth]{example-image-b}
        \subcaption{Variante B}
        \label{fig:nebeneinander-b}
    \end{minipage}
    \caption{Zwei Bilder nebeneinander}
    \label{fig:nebeneinander}
\end{figure}

\section{Drei Bilder in einer Zeile}

Drei Bilder gleichmäßig verteilt:

\begin{figure}[htbp]
    \centering
    \begin{minipage}{0.3\textwidth}
        \centering
        \includegraphics[width=\textwidth]{example-image-a}
        \subcaption{Bild 1}
    \end{minipage}
    \hfill
    \begin{minipage}{0.3\textwidth}
        \centering
        \includegraphics[width=\textwidth]{example-image-b}
        \subcaption{Bild 2}
    \end{minipage}
    \hfill
    \begin{minipage}{0.3\textwidth}
        \centering
        \includegraphics[width=\textwidth]{example-image-c}
        \subcaption{Bild 3}
    \end{minipage}
    \caption{Drei Bilder nebeneinander}
    \label{fig:drei-nebeneinander}
\end{figure}

\section{Von Text umflossenes Bild}
\index{wrapfig (Paket)}

Mit dem Paket \texttt{wrapfig} lässt sich Text um ein Bild fließen:

\begin{wrapfigure}{r}{0.4\textwidth}
    \centering
    \includegraphics[width=0.38\textwidth]{example-image}
    \caption{Von Text umflossen}
    \label{fig:wrapfig}
\end{wrapfigure}

Hier steht nun Fließtext, der das Bild auf der rechten Seite umfließt.
Dies ist nützlich, wenn man Platz sparen möchte und das Bild nicht die
ganze Breite einnehmen soll. Der Text läuft automatisch um das Bild herum
und setzt sich danach wieder über die gesamte Breite fort.

Weitere Zeilen Text, um den Umfluss-Effekt zu demonstrieren.
\LaTeX\ sorgt dafür, dass der Abstand zwischen Text und Bild angenehm
ist und kein Text über das Bild ragt.

% ============================================================
\chapter{Float-Positionierung}
\label{sec:float}
\index{Float-Positionierung}
% ============================================================

\section{Positionierungsangaben}

\LaTeX\ bietet verschiedene Optionen für die Platzierung von Gleitobjekten\index{Gleitobjekt} \cite{lamport1994}:

\begin{description}
    \item[\texttt{h}] \textbf{h}ere -- möglichst an der aktuellen Stelle
    \item[\texttt{t}] \textbf{t}op -- oben auf der Seite
    \item[\texttt{b}] \textbf{b}ottom -- unten auf der Seite
    \item[\texttt{p}] floating \textbf{p}age -- auf einer separaten Gleitseite
    \item[\texttt{H}] \textbf{H}ere -- genau hier (erfordert Paket \texttt{float})
\end{description}

\section{Bild erzwingen (H)}

\begin{figure}[H]  % H = genau hier
    \centering
    \includegraphics[width=0.4\textwidth]{example-image}
    \caption{Dieses Bild wird genau hier eingefügt}
    \label{fig:erzwungen}
\end{figure}

% ============================================================
\chapter{Referenzen auf Abbildungen}
\label{sec:referenzen}
\index{Referenz!Abbildung}
% ============================================================

Abbildungen lassen sich mit \texttt{\textbackslash label}\index{label} und
\texttt{\textbackslash ref}\index{ref} referenzieren:

\begin{itemize}
    \item Abbildung~\ref{fig:breite60}: Bild mit 60\,\% Textbreite
    \item Abbildung~\ref{fig:gedreht}: Gedrehtes Bild
    \item Abbildung~\ref{fig:nebeneinander}: Zwei Bilder nebeneinander
    \item Abbildung~\ref{fig:wrapfig}: Von Text umflossenes Bild
\end{itemize}